
\documentclass[9pt,a4paper,twocolumn,twoside]{rho-class/rho}
\usepackage[english]{babel}
\usepackage[backend=biber,style=ieee]{biblatex}
\addbibresource{bibliography.bib}

%----------------------------------------------------------
% Title
%----------------------------------------------------------

\doctype{Fintech Course Group Project}
\title{Data-driven customer profiling for financial institutions}

%----------------------------------------------------------
% Authors
%----------------------------------------------------------

\author[a]{Simone Zani}
\author[a]{Giulia Talà}
\author[a]{Tommaso Alessio Baresi}
\author[a]{Alberto Toia}
\author[a]{Marco Amarilli}

\affil[a]{Politecnico di Milano}

%----------------------------------------------------------

\dates{\today}

\corres{Course: Fintech}

\docinfo{Clustering Analysis on Customer Data}

\journalname{Politecnico di Milano}
\journal{Fintech}
\theday{\today}
\vol{}
\no{}

%----------------------------------------------------------
% Abstract (vuoto perché nel tuo report non c’era)
%----------------------------------------------------------

\begin{abstract}
Customer segmentation is an important task for financial institutions, especially when dealing with heterogeneous data. In this project, we focus on clustering bank customers using a dataset that includes both numerical and categorical variables.

To handle mixed-type data, we define a distance measure that combines numerical and categorical components. We test different distance metrics and optimize their relative weight using a parameter $\alpha$. Clustering is performed with the K-Medoids algorithm and evaluated using the silhouette score. We also validate the results using UMAP-based dimensionality reduction and stability analysis.

The results show that giving more importance to categorical variables improves clustering performance. The best configuration achieves a silhouette score of 0.6788 and identifies 10 clusters. The UMAP validation confirms that the clustering is stable across different runs.

Furthermore, we extend the behavioral segmentation by assigning a tangible economic value to each cluster. By deploying Logistic Classifiers for retail product propensities and leveraging European and Italian market margins alongside channel-specific service costs, we compute the expected net revenue. This allows us to translate data-driven personas into actionable marketing strategies.

Overall, the analysis suggests that properly combining different types of variables is important for obtaining meaningful customer segmentation, and projecting these profiles onto macroeconomic constraints produces highly actionable results.
\end{abstract}

\keywords{Customer Segmentation, Clustering, Mixed Data, K-Medoids, UMAP}

%----------------------------------------------------------

\begin{document}

\maketitle
\thispagestyle{firststyle}

%----------------------------------------------------------

\section{Introduction}

In this work, we aim to segment bank customers into homogeneous groups based on the information contained in their feature vectors. The dataset describes clients through a combination of demographic, socio-economic, and behavioral variables, including both numerical features (e.g., income, wealth, financial activity) and categorical features (e.g., gender, job, geographical area).

The objective of the analysis is to identify meaningful customer personas that summarize different financial profiles and behaviors within the population.

A challenge in this setting is the presence of mixed-type data. Standard clustering techniques (e.g. K-Means \cite{Ikotun2023}) are designed for purely numerical variables and are not directly applicable in this context. For this reason, we go beyond and focus on defining an appropriate distance measure that can consistently combine numerical and categorical information.

We then use this distance to perform clustering and evaluate how different modeling choices affect the quality, stability, and interpretability of the resulting segmentation.

Finally, we extend the behavioral profiling by assigning a tangible economic value to the derived segments. Through a combination of supervised numerical regression (Gradient Boosting and Logistic models) and categorical profiling via a Generalized Learning Vector Quantization (GLVQ) classification network, we cross-validate the structural integrity of the clusters' profitability, translating raw patterns into actionable revenue predictors.

%----------------------------------------------------------

\section{Methodology}

\subsection{Data Preprocessing}
We started from a dataset of 5000 bank clients and we designed the preprocessing phase in order to improve both data consistency and comparability across features.

First, we imputed missing values using the median for numerical variables and the mode for categorical ones, then we also applied a domain-based consistency filter to remove unrealistic cases such as underaged individuals with working occupations. In addition we detected outliers using an Isolation Forest model with contamination level fixed at 1\% leading to a final sample consisting of 4950 observations.

Finally, we rescaled numerical variables to the $[0,1]$ interval through Min-Max normalization, while for categorical variables we used label encoding for Hamming-based distances and one-hot encoding for Tanimoto-based distances.

\subsection{Distance Construction}
Since the dataset contains both numerical and categorical variables, we carefully define a distance measure that can combine the two.

For numerical variables, we consider three metrics: Manhattan (L1), Euclidean (L2), and Canberra distance \cite{Lance1967}.

For categorical variables, we use two approaches. We compute Hamming distance \cite{Hamming1950} on label-encoded variables and Tanimoto distance \cite{Tanimoto1958} on one-hot encoded variables.

We combine the two components using a weighted distance $d$:
\newline
\[
d(x,y) = \alpha d_{\mathrm{num}}(x,y) + (1-\alpha)d_{\mathrm{cat}}(x,y),
\]
where $\alpha \in [0,1]$ controls the relative importance of numerical and categorical variables.

Gower distance \cite{Gower1971} is taken into account as well, with the advantage of defining a single distance for both numerical and categorical variables.

\subsection{Basic Clustering (Neutral Alpha)}

We start by performing clustering using a baseline distance configuration. In this setting, we apply the K-Medoids algorithm \cite{Kaufman2008}, which is well suited for mixed-type data as it operates directly on distance matrices and does not rely on Euclidean assumptions.

The number of clusters is explored in the range $k=2,\dots,10$, and each solution is evaluated using the silhouette score \cite{Rousseeuw1987} as the main criterion. To ensure robustness, we repeat the clustering procedure with multiple random seeds and compute the average performance.

As a first step, we define the distance using a neutral value of $\alpha$, set proportionally to the number of features in the numerical and categorical blocks. This corresponds to assigning equal weight to each variable and is consistent with the weighting scheme used by Gower distance \cite{Gower1971}, which we adopt as a benchmark.

This baseline allows us to compare different distance specifications under a common and unbiased setting. 

\subsection{Weighted Clustering}

After establishing the baseline, we introduce a weighted distance to better control the contribution of numerical and categorical variables. Instead of fixing $\alpha$, we perform a grid search over multiple values and select the optimal configuration based on clustering performance.

For each combination of numerical and categorical distance, we evaluate different values of $\alpha$ and identify the best pair $(\alpha, k)$ according to the average silhouette score \cite{Rousseeuw1987} across multiple runs. This approach allows each distance specification to be compared at its optimal setting, avoiding biases introduced by a fixed weighting scheme.

\begin{figure}[H]
    \centering
    \includegraphics[width=\linewidth]{figures/alpha_x_silh.png}
    \caption{Silhouette score trend as a function of the $\alpha$ parameter.}
    \label{fig:alpha_silh}
\end{figure}

\begin{figure*}[ht]
    \centering
    \begin{subfigure}[b]{0.48\textwidth}
        \centering
        \includegraphics[width=\linewidth]{figures/umap2d.png}
        \caption{2D UMAP projection (Weighted results)}
    \end{subfigure}
    \hfill
    \begin{subfigure}[b]{0.48\textwidth}
        \centering
        \includegraphics[width=\linewidth]{figures/umap3d.png}
        \caption{3D UMAP projection (Weighted results)}
    \end{subfigure}
    \caption{Dimensionality reduction via UMAP confirming cluster separation and stability.}
    \label{fig:umap_plots}
\end{figure*}

\begin{figure}[H]
    \centering
    \includegraphics[width=\linewidth]{figures/silhuette.png}
    \caption{Silhouette analysis for the optimal $k=10$ clusters.}
    \label{fig:silhouette_plot}
\end{figure}

\begin{figure*}[ht]
    \centering
    \includegraphics[width=0.8\textwidth]{figures/radar.png}
    \caption{Radar chart showing the average features across the identified customer segments.}
    \label{fig:radar_chart}
\end{figure*}

\begin{figure*}[ht]
    \centering
    \includegraphics[width=0.9\textwidth]{figures/key.png}
    \caption{Key feature distribution broken down by clusters.}
    \label{fig:key_features}
    \end{figure*}

\subsection{UMAP-Based Validation}

To further validate the robustness of the results, we adopt an alternative clustering approach based on dimensionality reduction. We first apply UMAP \cite{McInnes2018} to project the data into a lower-dimensional space, and then perform clustering using K-Medoids \cite{Kaufman2008} on the resulting representation.

We test multiple embedding dimensionalities (8, 12, and 15) and evaluate both clustering quality and stability. Stability is assessed by repeating the pipeline with different random seeds and computing the Adjusted Rand Index \cite{Hubert1985} between the resulting partitions.

\subsection{Economic Valuation and Propensity Modeling}

To bridge the gap between behavioral clustering and tangible macroeconomic profitability, we develop a Customer Lifetime Value (CLV) proxy. 
Instead of assuming static engagement rates across the bank, we train Logistic Regression Classifiers to predict the propensity of each client to purchase or hold specific retail products (Savings, Investments, ESG Funds, and Mortgages). These target models use sociodemographic and financial indicators as features, outputting a probability score $[0, 1]$ for each product. The classifiers proved highly suited for this task due to their interpretability.

With the propensities computed, we estimate the Expected Annual Gross Revenue by multiplying the purchase probability by the average ticket size of that product category (e.g., €50,000 for Savings, €150,000 for Mortgages) and the respective Italian market margins. We anchor these margins to current European/Italian macroeconomic rates, setting the Net Interest Margin for Savings at 1.8\%, and the retroactive distribution fees for Investments, ESG Funds, and Mortgages at 1.2\%, 1.4\%, and 1.5\%, respectively.

$$ \text{Gross Revenue} = \sum_{p \in \text{Products}} (\text{Prop}_p \times \text{AUM Proxy}_p \times \text{Margin}_p) $$

Finally, to compute the Net Revenue (CLV proxy), we adjust the gross figures by accounting for the specific Cost to Serve (CTS). We establish a baseline physical service cost of €300 per client annually, which is scaled down based on the digital adoption of the cluster. Digital-native clusters receive a 0.70x multiplier reflecting structural IT scaling efficiencies (bypassing expensive physical advisory), hybrid segments a 0.85x multiplier, and branch-dependent traditional clusters retain the full 1.0x physical cost burden.

%----------------------------------------------------------

\section{Results}

\subsection{{Clustering Results and Distance Comparison}}

Under the neutral (non-weighted) configuration, clustering performance is moderate. Tanimoto-based distances outperform Hamming-based ones, identifying $k=10$ clusters with silhouette scores around 0.40–0.46, while Hamming and Gower configurations remain much weaker, with silhouette values close to 0.16.

When allowing for weighted distances, all configurations achieve their best performance at $\alpha = 0.2$, meaning that categorical variables receive 80\% of the total weight.

This results in a substantial improvement in clustering quality. The best configuration, based on Manhattan distance and Tanimoto encoding, reaches a silhouette score of 0.6788, compared to approximately 0.16 obtained with Gower distance.

\begin{table}[H]
\centering
\caption{Best clustering results with optimized alpha}
\begin{tabular}{lccc}
\toprule
Distance & Alpha & k & Silhouette \\
\midrule
L1 + Tanimoto       & 0.2 & 10 & 0.6788 \\
Canberra + Tanimoto & 0.2 & 10 & 0.6671 \\
L2 + Tanimoto       & 0.2 & 10 & 0.6579 \\
L1 + Hamming        & 0.2 & 10 & 0.6492 \\
\bottomrule
\end{tabular}
\end{table}

We also observe that the improvement is consistent across all configurations. This suggests that the gain is not driven by a specific metric choice, but by a better balance between numerical and categorical information.

\subsection{UMAP Validation}

We then validate the results using UMAP-based clustering. All configurations consistently select $k=10$ clusters, confirming the robustness of the segmentation.

The best configuration (8-dimensional embedding) achieves a silhouette score of 0.7813.

\begin{table}[H]
\centering
\caption{UMAP clustering performance}
\begin{tabular}{lccc}
\toprule
Configuration & k & Silhouette & ARI \\
\midrule
UMAP 8D  & 10 & 0.7813 & 0.948 \\
UMAP 12D & 10 & 0.7577 & 0.933 \\
UMAP 15D & 10 & 0.7593 & 0.954 \\
\bottomrule
\end{tabular}
\end{table}

We also evaluate stability across multiple runs. All configurations achieve ARI values above 0.93, indicating very high consistency across different initializations. Among them, the 8-dimensional embedding provides the best trade-off between separation and stability.

The consistency between distance-based clustering and UMAP-based clustering suggests that the structure we identify is intrinsic to the data rather than based on a specific modeling choice.

\subsection{Cluster Interpretation}

The final segmentation consists of 10 clusters with well-balanced sizes, without degenerate or extremely small groups. By analyzing the demographic, financial, and behavioral profiles of the centroids, we can interpret these clusters as distinct customer personas. Categorical variables such as Job, Gender, and Geographical Area play a central role in separating the groups, confirming the importance of our weighted distance construction.

The 10 clusters naturally pair into five demographic profiles, often bridging gender divides within similar socio-economic strata:

\begin{itemize}
    \item \textbf{Clusters 5 \& 6 (Mass-Market Employees - North):} Representing over half of the portfolio, these are middle-aged professionals (mid-50s) located in Northern Italy. Cluster 5 (female) is the single largest cluster, while Cluster 6 (male) shows slightly higher income, wealth, and debt levels. Both exhibit strong digital activity, solid income, and high ESG affinity, making them the prime targets for digital and sustainable products.
    \item \textbf{Clusters 0 \& 2 (Central-Region Employees):} These segments represent employed individuals in their mid-to-late 50s based in Central Italy. Cluster 0 (male) and Cluster 2 (female) share solid savings rates and stable financial profiles. While similar in income to the Northern employees, their regional allocation and lower debt levels define a more traditional, ``active saver'' profile.
    \item \textbf{Clusters 3 \& 7 (Retirees - North):} Comprising retired individuals (average age mid-to-late 70s) in the North. Cluster 3 characterizes retired women, and Cluster 7 retired men. They are defined by low debt and very low technological engagement but exhibit significant ESG awareness. Unsurprisingly, they represent the bank's core demographic for traditional human-assisted channels and capital preservation.
    \item \textbf{Clusters 1 \& 4 (Pre-Retirement Unemployed):} These clusters capture individuals in a transitional financial phase (early-to-mid 60s, unemployed) in the North. Cluster 1 (female) and Cluster 4 (male) show middle-tier income but notably low digital engagement and reduced financial activity. Their profiles suggest intact income (likely early retirement or redundancy provisions) but a general drop in active banking interactions.
    \item \textbf{Clusters 8 \& 9 (Southern Family Employees):} Geographically distinct, these employed clients in the South (late 50s) represent the largest family units in the dataset (averaging nearly 3 members per household). Cluster 8 portrays a male ``family man'' profile with moderate debt, while Cluster 9 consists of female clients showing significant real asset wealth. Both clusters exhibit an exceptionally strong fit for life-event banking and family-oriented credit products.
\end{itemize}

Overall, we observe a clear structure where life stage (youth/employment vs. retirement), geography, and family size dictate the financial needs and digital behaviors of each group, fully validating our dual continuous-categorical distance metric.

\subsection{Strategic Economic Value and Marketing Insights}

The synthesis of behavioral clustering and economic valuation reveals five distinct strategic macro-families within our 10 unsupervised clusters. By superimposing the economic value scores over the demographic profiles, we can isolate highly targeted marketing strategies for each segment, thoroughly summarized in Table \ref{tab:strategy_summary} and visually grouped in Figure \ref{fig:net_revenue}.

Specifically, \textbf{Volume Engines} and the \textbf{Stable Core} represent the core profitability drivers through highly digital, mass-affluent cross-selling and wealth accumulation. \textbf{Family Borrowers} dominate mortgage propensities, representing key life-event lending targets. On the lower revenue spectrum, \textbf{Transitional} clients and \textbf{ESG Retirees} necessitate trust retention and capital preservation strategies, owing to their inherently higher physical channel costs and different lifecycle stages.

\begin{table}[htbp]
\centering
\caption{Summary of Strategic Macro-Families and Economic Valuation}
\resizebox{\columnwidth}{!}{%
\begin{tabular}{llccl}
\toprule
\textbf{Strategic Family} & \textbf{Clusters} & \textbf{Total Margin} & \textbf{Per Capita} & \textbf{Key Products} \\
\midrule
Volume Engines & 5 \& 6 & $\sim$ \texteuro 4.48M & $\sim$ \texteuro 1,750 & Savings, Mortgages, ESG \\
Stable Core & 0 \& 2 & $\sim$ \texteuro 0.92M & $\sim$ \texteuro 1,630 & Savings, Investments \\
Family Borrowers & 8 \& 9 & $\sim$ \texteuro 0.80M & $\sim$ \texteuro 1,750 & Mortgages (Dominant) \\
Transitional & 1 \& 4 & $\sim$ \texteuro 0.58M & $\sim$ \texteuro 1,100 & Low generalized \\
ESG Retirees & 3 \& 7 & $\sim$ \texteuro 0.48M & $\sim$ \texteuro 580 & ESG (Dominant) \\
\bottomrule
\end{tabular}%
}
\label{tab:strategy_summary}
\end{table}

\begin{figure*}[ht]
    \centering
    \includegraphics[width=0.8\textwidth]{figures/net_revenue_clusters.png}
    \caption{Strategic macro-families grouped by Expected Annual Net Revenue.}
    \label{fig:net_revenue}
\end{figure*}

%----------------------------------------------------------

\section{Conclusion}

In this project, we show that customer segmentation can be significantly improved by carefully combining numerical and categorical information.

The best-performing model uses Manhattan distance and Tanimoto encoding with $\alpha = 0.2$, producing a stable segmentation with 10 clusters. The UMAP-based validation confirms the robustness of the result.

Overall, we find that categorical variables play a dominant role in explaining customer heterogeneity, and properly weighting them is essential for obtaining meaningful segmentation.

Furthermore, moving beyond standard unsupervised profiling, we prove that integrating specific market-margin multipliers and channel-based cost-to-serve metrics yields clear, actionable measures of Expected Net Revenue (CLV proxy). Categorical and numerical variables are crucial for identifying behavioral profiles, but superimposing predictive product propensities empowers institutions to map these personas into precise economic value drivers and strategic macro-families.

In conclusion, properly weighting heterogenous variables along with tangible financial constraints enables the end-to-end transformation of raw data into a meaningful, strategy-oriented segmentation.

\nocite{*}
\printbibliography

%----------------------------------------------------------

\end{document}